\documentclass[12pt,a4paper]{article}
% ============================================================
%  Оформление отчёта. Обычно здесь ничего менять не нужно.
%  Собирается pdfLaTeX (по умолчанию, в т. ч. в Overleaf),
%  XeLaTeX или LuaLaTeX.
% ============================================================

\usepackage{iftex}
\ifPDFTeX
  \usepackage[T2A]{fontenc}
  \usepackage[utf8]{inputenc}
  \usepackage{tempora}                 % Times-подобный шрифт с кириллицей
\else
  \usepackage{fontspec}
  \IfFontExistsTF{Times New Roman}
    {\setmainfont{Times New Roman}}
    {\setmainfont{CMU Serif}}
  \IfFontExistsTF{Consolas}
    {\setmonofont{Consolas}[Scale=MatchLowercase]}
    {\IfFontExistsTF{Courier New}
      {\setmonofont{Courier New}}
      {\setmonofont{CMU Typewriter Text}}}
\fi
\usepackage[english,russian]{babel}

% --- Поля и абзацы (как в исходном doc-шаблоне курса) ---
\usepackage[a4paper, margin=2.54cm]{geometry}
\setlength{\parindent}{0pt}
\setlength{\parskip}{4pt plus 1pt}
\frenchspacing

% --- Заголовки разделов без номеров ---
\usepackage{titlesec}
\setcounter{secnumdepth}{0}
\titleformat{\section}{\large\bfseries}{}{0pt}{}
\titleformat{\subsection}{\normalsize\bfseries}{}{0pt}{}
\titlespacing*{\section}{0pt}{14pt plus 4pt minus 2pt}{4pt plus 2pt}
\titlespacing*{\subsection}{0pt}{10pt plus 2pt minus 2pt}{3pt plus 1pt}

% --- Таблицы, рисунки, подписи ---
\usepackage[table]{xcolor}
\usepackage{graphicx}
\usepackage{tabularx}
\usepackage{float}
\usepackage{caption}
\DeclareCaptionLabelSeparator{spacedendash}{~-- }
\captionsetup{labelsep=spacedendash, font=small, singlelinecheck=false, justification=raggedright}
\captionsetup[figure]{justification=centering, singlelinecheck=true}
\captionsetup[table]{position=top, skip=4pt}
\DeclareCaptionType[fileext=lol]{listing}[Листинг][Листинги]
\captionsetup[listing]{position=top, skip=4pt, hypcap=false}
\addto\captionsrussian{%
  \renewcommand{\figurename}{Рисунок}%
  \renewcommand{\tablename}{Таблица}%
}
\newcommand{\thead}[1]{\multicolumn{1}{|c|}{\cellcolor{gray!20}\textbf{#1}}}
% Колонки таблиц без растягивания пробелов: L – резиновая, P{ширина} – фиксированная
\newcolumntype{L}{>{\raggedright\arraybackslash}X}
\newcolumntype{P}[1]{>{\raggedright\arraybackslash}p{#1}}

% --- Код и вывод терминала ---
% code     – фрагмент исходного кода;
% terminal – вывод программы, strace и т. п.;
% \inputterminal{файл} / \inputcode{файл} – то же, но из файла.
\usepackage{fvextra}
\fvset{breaklines=true, breakanywhere=true,
       frame=leftline, framerule=1.5pt, framesep=6pt,
       rulecolor=\color{gray!60}, xleftmargin=4pt}
\DefineVerbatimEnvironment{code}{Verbatim}{fontsize=\small}
\DefineVerbatimEnvironment{terminal}{Verbatim}{fontsize=\footnotesize}
\newcommand{\inputcode}[2][]{\VerbatimInput[fontsize=\small,#1]{#2}}
\newcommand{\inputterminal}[2][]{\VerbatimInput[fontsize=\footnotesize,#1]{#2}}
% Листинг с подписью сверху:
%   \begin{listingblock}{Подпись} \begin{code}...\end{code} \end{listingblock}
\newenvironment{listingblock}[1]
  {\par\medskip\captionof{listing}{#1}\nopagebreak\vspace{-\topsep}\vspace{-\parskip}}
  {\par}

% Приложение: с новой страницы, заголовок «Приложение А. ...»,
% листинги внутри нумеруются А.1, А.2, ...
\newcounter{appendixno}
\newcommand{\appendixsection}[2]{%
  \clearpage
  \section{Приложение #1. #2}%
  \stepcounter{appendixno}%
  \setcounter{listing}{0}%
  \renewcommand{\thelisting}{#1.\arabic{listing}}%
  \renewcommand{\theHlisting}{app\arabic{appendixno}.\arabic{listing}}}

% --- Ссылки ---
\usepackage[hidelinks]{hyperref}
\usepackage{xurl}

% --- Подсказки шаблона: серый курсив, удалить перед сдачей ---
\newcommand{\hint}[1]{\leavevmode{\color{gray}\itshape #1}}
% Рамка-заглушка вместо рисунка
\newcommand{\placeholder}[2][4.5cm]{%
  \fbox{\parbox[c][#1][c]{0.95\linewidth}{\centering\hint{#2}}}}

\newcommand{\ReportStudent}{В. Е. Авраменко}
\newcommand{\ReportGroup}{М8О-204БВ-25}
\newcommand{\ReportTeacher}{Н. А. Гапонов}
\newcommand{\ReportDate}{18.11.2026}
\newcommand{\ReportYear}{2026}
\newcommand{\ReportCourse}{Операционные системы и архитектура компьютера}
\begin{document}
% Титульный лист
\begin{titlepage}
  \centering\large

  МОСКОВСКИЙ АВИАЦИОННЫЙ ИНСТИТУТ\\
  (НАЦИОНАЛЬНЫЙ ИССЛЕДОВАТЕЛЬСКИЙ УНИВЕРСИТЕТ)\par
  \vspace{18pt}
  Институт №8 «Компьютерные науки»\\
  Кафедра 806 «Вычислительная математика и программирование»\par

  \vspace{\stretch{1.2}}

  {\bfseries Отчёт по лабораторной работе №5\\
  Вариант 9\\
  по курсу «\ReportCourse»\par}

  \vspace{\stretch{1.6}}

  \begin{flushright}
    \begin{tabular}{r@{\hspace{0.6em}}l}
      Выполнил:      & \ReportStudent \\[2pt]
      Группа:        & \ReportGroup   \\[2pt]
      Преподаватель: & \ReportTeacher \\[2pt]
      Дата:          & \ReportDate    \\[2pt]
      Оценка:        & \rule[-2pt]{4.2cm}{0.4pt} \\
    \end{tabular}
  \end{flushright}

  \vspace{\stretch{1}}

  Москва, \ReportYear
\end{titlepage}
\setcounter{page}{2}

\section{Цель работы}
Научиться создавать динамические библиотеки и использовать их двумя способами: связывать исполняемый файл с библиотекой при сборке и явно загружать библиотеку во время работы программы.

\section{Описание решения}
Вариант 9 включает функцию 2 (производная $\cos x$) и функцию 3 (число простых на включительном отрезке $[A,B]$). Обе библиотеки экспортируют функции с одинаковыми сигнатурами.
\inputcode[firstline=4,lastline=5]{../src/contracts.h}

В первой реализации приближение производной равно $(\cos(A+h)-\cos(A))/h$. Для проверки числа $n$ перебираются делители от 2 до $n-1$ до первого делителя. Во второй реализации используется центральная разность $(\cos(A+h)-\cos(A-h))/(2h)$ и сегментированное решето Эратосфена. Решето сначала находит простые до $\lfloor\sqrt{B}\rfloor$. Затем оно вычёркивает их кратные в блоках отрезка, начиная с первого кратного в блоке, которое не меньше квадрата простого. Размер временного блока постоянен.

Обе программы используют один обработчик строк. Команда \texttt{1 A deltaX} принимает конечные числа и ненулевой шаг, угол задаётся в радианах. Команда \texttt{2 A B} принимает натуральные границы, где $A\leq B$. Ошибка ввода сопровождается сообщением в stderr. Стоимость наивной проверки растёт с величиной чисел, поэтому широкие интервалы удобнее проверять после переключения на решето.

\texttt{linked} связывается с \texttt{libforward.so} при сборке. Компоновщик записывает зависимость в ELF, и загрузчик обрабатывает её при запуске. \texttt{runtime} открывает первую библиотеку через \texttt{dlopen} и переключает её командой \texttt{0}. \texttt{dlsym} получает адреса функций, а \texttt{dlerror} проверяет ошибку поиска. При ошибке загрузки новой библиотеки или поиска символа программа сообщает причину и сохраняет прежнюю реализацию. При выходе программа закрывает текущий дескриптор.

\begin{listingblock}{Получение адресов функций и смена реализации в runtime.c}
\inputcode[firstline=48,lastline=61,fontsize=\footnotesize]{../src/runtime.c}
\end{listingblock}

\clearpage
\section{Результаты работы}
Для $A=0$ и $h=0{,}5$ истинное значение производной равно нулю. После переключения результат меняется с $-0{,}24483487$ на нуль. Обе реализации насчитывают 25 простых от 1 до 100. Возврат к первой библиотеке восстанавливает исходный результат производной.
\inputterminal{inc/demo.txt}

Проверка сравнила обе реализации подсчёта с независимой проверкой делителей на 55 отрезках. Для решета отдельно подтверждены 9592 простых от 1 до 100000 и принадлежность числа 2147483647 к простым. Также проверены неверный ввод, отсутствие библиотеки, отсутствие символа в библиотеке и ошибка записи в stdout. Прогоны двух программ под AddressSanitizer и UndefinedBehaviorSanitizer завершились без диагностических сообщений.

\clearpage
\section{Проверка загрузки библиотек через strace}
Для повторения опыта нужно собрать программы командой \texttt{make} и запустить следующие команды из корня проекта. Ключ \texttt{-o} сохраняет трассу отдельно от вывода программы, \texttt{-yy} подписывает файловые дескрипторы, а \texttt{-s 128} позволяет видеть команды целиком.
\begin{terminal}
syscalls=execve,openat,mmap,mprotect,close,munmap,read,write,exit_group
printf '1 0 0.5\n2 1 100\n' |
  strace -yy -s 128 -e trace="$syscalls" -o build/linked-full.trace ./build/linked
printf '1 0 0.5\n0\n1 0 0.5\n2 1 100\n0\n1 0 0.5\n' |
  strace -yy -s 128 -e trace="$syscalls" -o build/runtime-full.trace ./build/runtime
\end{terminal}
В выдержках ниже \texttt{\$LAB} обозначает абсолютный путь к проекту. Сохранены успешные открытия библиотек варианта, начало отображения каждой библиотеки, отображение её исполняемого сегмента, закрытие файла и соответствующие вызовы \texttt{munmap}. Системные зависимости, другие сегменты и промежуточные проверки путей видны в полных файлах трассировки. Адреса и номера дескрипторов могут отличаться при повторном запуске.

\subsection{Программа 1. Библиотека из зависимостей ELF}
\texttt{readelf} показывает \texttt{DT\_NEEDED} для \texttt{libforward.so} в \texttt{linked}. \texttt{RUNPATH} со значением \texttt{\$ORIGIN} задаёт поиск библиотеки рядом с исполняемым файлом. У \texttt{runtime} зависимости от библиотек варианта в ELF отсутствуют.
\inputterminal{inc/elf.txt}

\inputterminal{inc/trace-linked.txt}
Сразу после \texttt{execve} динамический загрузчик обрабатывает зависимости: \texttt{openat} возвращает файловый дескриптор 3, \texttt{mmap} отображает библиотеку в адресное пространство процесса, а \texttt{PROT\_EXEC} разрешает исполнение кода в соответствующем сегменте. \texttt{close(3)} закрывает файл, при этом отображённые страницы остаются доступны. Первый \texttt{read(0, ...)} начинается уже после загрузки. Вызовы контрактов выполняются в пользовательском пространстве, а \texttt{write(1, ...)} печатает их результаты. При завершении процесса ядро освобождает его адресное пространство.

\clearpage
\subsection{Программа 2. Начальная явная загрузка и первая команда 0}
\inputterminal[firstline=1,lastline=15]{inc/trace-runtime.txt}
Программа вызывает \texttt{dlopen} для \texttt{./build/libforward.so} в начале \texttt{main}. В трассе относительный путь передаётся в \texttt{openat} с \texttt{AT\_FDCWD}, поэтому отсчёт идёт от текущего каталога. Флаг \texttt{RTLD\_NOW} требует разрешить неопределённые символы загружаемой библиотеки до возврата из \texttt{dlopen}. Флаг \texttt{RTLD\_LOCAL} исключает её символы из области поиска зависимостей для последующих загрузок. Через полученный дескриптор их можно найти с помощью \texttt{dlsym}. Загрузчик закрывает файловый дескриптор после отображения. Программа хранит дескриптор загрузчика, полученный от \texttt{dlopen}, до вызова \texttt{dlclose}.

Один вызов \texttt{read(0, ...)} принял все строки из канала. \texttt{getline} разбирает их по одной из буфера. Сначала программа печатает результат правой разности, затем обрабатывает первую команду \texttt{0}. Следующий \texttt{openat} открывает \texttt{libcentral.so}, а \texttt{mmap} отображает её код. После поиска двух символов программа вызывает \texttt{dlclose} для прежней библиотеки. В этом запуске ему соответствует \texttt{munmap} с адресом начала отображения \texttt{libforward.so} (\texttt{0x78ad305e3000}) и длиной 16408 байт. Затем \texttt{write} подтверждает смену библиотеки. Следующие результаты равны 0 для центральной разности и 25 для решета.

\clearpage
\subsection{Возврат к первой библиотеке и завершение}
\inputterminal[firstline=16]{inc/trace-runtime.txt}
Вторая команда \texttt{0} заново открывает \texttt{libforward.so}. В этом прогоне загрузчик повторно использовал тот же виртуальный адрес. Вызов \texttt{munmap} с адресом \texttt{0x78ad305de000} и длиной 16432 байта снимает отображение \texttt{libcentral.so}. После EOF, который виден по возврату нуля из \texttt{read(0, ...)}, программа освобождает последнее отображение \texttt{libforward.so}. \texttt{exit\_group(0)} завершает программу со статусом 0.

Вызовы \texttt{dlopen}, \texttt{dlsym} и \texttt{dlclose} обрабатывает загрузчик в пользовательском пространстве. \texttt{strace} показывает его обращения к ядру. Поиск \texttt{dlsym} идёт по уже отображённым таблицам ELF и в трассе отдельной строкой не появляется. \texttt{dlclose} уменьшает счётчик ссылок на библиотеку. В этом опыте после закрытия последней ссылки её отображение снимается.

\section{Выводы}
При связывании во время сборки программе достаточно обычных вызовов функций, а наличие библиотеки проверяется загрузчиком при старте. При явной загрузке можно выбирать реализацию во время работы без пересборки программы. Для этого программа самостоятельно обрабатывает ошибки загрузки и поиска символов и управляет временем жизни дескриптора. При одинаковых входах реализация подсчёта не меняет результат, а разные формулы конечных разностей дают различное приближение производной.
\end{document}
